\section{From the boundary value to the circuit lower bound}\label{sec:border}

\begin{proof}[Proof of Theorem~\ref{thm:stretch}]
Fix a characteristic admitting ports and the corresponding compressor
identities and clocks. Suppose that at a sufficiently large $n$ the
coefficient vector of $P_n^*$ lies in the degree-component border of
circuits of size $m_0=\lfloor 2^{n^\sigma}\rfloor$. In positive
characteristic use the given prime $p$. In characteristic zero apply
Lemma~\ref{lem:descent-prime} to the parse table. It preserves the
image-closure inclusion and the fixed rational compressor identities
at a prime $p$ of polynomial bit length in $M$. This prime may
depend on $n$. In both cases we evaluate compressor constants in
$\mathbb F_p$.

Apply Lemma~\ref{lem:pins} on the curve of
Proposition~\ref{prop:descent-curve} with $L=R(2n-1)$. Write
$W_\mu=W_{1,\mu}$, use the padded binary indexing, and choose an
index $H_*$ at which $\Omega(H_*)\ne0$. All matrices, their entry
procedure and its advice are now fixed. Include $H_*$, the advice,
the size parameters and the public program in a description $D(H)$,
reserving one field for a variable index $H$. Define
\[
 \begin{aligned}
 Z(H)&=Z(D(H))\bmod p, & A_H&=A_{D(H)},\\
 W(H)&=\sum_{|\mu|\le L}([X^\mu]A_H)\,W_\mu .
 \end{aligned}
\]
For any choice of public program, $A_H=F_{D(H)}(Y(X))$ is regular on the
chosen branch and has value $F_{D(H)}(c)=Z(H)$. By linearity in the
coefficient in Lemma~\ref{lem:pins},
\begin{equation}\label{eq:border-row}
                   W(H)\Omega=Z(H)\Omega\qquad\text{for every }H .
\end{equation}
In particular, row $H$ of this identity computes $Z(H)\Omega(H)$.

Choose the public program on input $D(H)$ to count, modulo $p$,
\begin{equation}\label{eq:border-column}
  \mathbf 1_{\{H=H_*\}}+\sum_I W(I)_{I,H}
  =\mathbf 1_{\{H=H_*\}}+
       \sum_{I,\mu}([X^\mu]A_I)(W_\mu)_{I,H}.
\end{equation}
One branch accepts precisely when $H=H_*$. The other guesses $I$ and
a valid exponent vector $\mu$ in their canonical binary encodings and
replaces $H$ by $I$ in a copy of the description. Supply the active
words of $D(I)$ and the exponent vector $\mu$ as one input of at most
$R$ words, with the unused words of $D(I)$ understood as zero.
Lemma~\ref{lem:substitution} counts $[X^\mu]A_I$ from this input at
the larger clock. Simulate it
at the smaller clock with private jobs by Theorem~\ref{thm:gap}(2).
On each completed path, evaluate its tallied compressor monomial in
$\mathbb F_p$. A further binary guess smaller than its representative
in $[0,p-1]$ gives exactly that many successes. Independently, run the
entry procedure from Lemma~\ref{lem:pins} for $(W_\mu)_{I,H}$, requiring
both procedures to accept. This multiplies their counts. Summing over
the guesses gives \eqref{eq:border-column}.

The fixed private machine determines the wiring generator before we
choose the public program. The larger-clock procedure uses $D(I)$ as
data for the pins and links of that wiring. The smaller-clock program
constructs this input by copying its stored instructions and advice,
changing only the field for the index. Thus the same code works at
every $H$. The fixed code and the $M^{O(1)}$ bits of data fit within
the reserved $R/10$ active words for large $n$. Apart from the
$o(T_-)$ simulation and the $O(R)$ cost of copying the description,
the prime-field computations and the entry path cost $M^{O(1)}=o(R)$
ticks. The entire program therefore fits $T_-$.

By Theorem~\ref{thm:gap}(1), the success count of this program equals
$Z(H)$ in $\mathbb F_p$. Multiply \eqref{eq:border-column} by
$\Omega(H)$ and sum over $H$. Interchanging the finite sums groups
row $I$ with $\Omega$ in \eqref{eq:border-row}, giving
\[
 \sum_H Z(H)\Omega(H)
 =\Omega(H_*)+\sum_I\bigl(W(I)\Omega\bigr)_I
 =\Omega(H_*)+\sum_I Z(I)\Omega(I).
\]
Cancellation contradicts the choice of $H_*$. Hence
$\bsz_F^{[n]}(P_n^*)>2^{n^\sigma}$ for all sufficiently large $n$
in the chosen characteristic.

To transfer the bound, use the signed verifier of
Section~\ref{sec:model}. It accepts the paths with output $+1$ or
$-1$, gives the negative-output bit weight $-1$, and attaches the
selected variables $Z_{i,x_i}$ to its output bits. Its length is
polynomial in $n$. Lemma~\ref{gadget:verifier} projects $\HC$ of order
at most $A n^B$ to $P_n^*$, and supplies the same projection for
$\Per$ outside characteristic two, for fixed $A,B>0$. At order $N$,
take $n=\lfloor(N/A)^{1/B}\rfloor$ and pad the graph by a forced path
or an identity block, respectively. As in Section~\ref{sec:comparison},
extracting the degree-$N$ component of a size-$s$ approximant before
projection costs $O((N+1)^2(s+1))$. The projection and the final
degree-$n$ extraction preserve the closure. Consequently this polynomial
overhead times $s+1$ must exceed $2^{n^\sigma}$. Taking any fixed
$0<\varepsilon<\sigma/B$ gives the asserted $2^{N^\varepsilon}$
lower bound for all sufficiently large $N$.

Keeping a distinct variable on each edge of the verifier graph gives
the port case of Corollary~\ref{cor:sparse-support}. Specializing those
variables to the verifier weights gives precisely the projection above,
and the padding preserves the degree bounds on the support.
\end{proof}

\begin{proof}[Completion of the proof of Theorem~\ref{thm:main}]
In each characteristic without ports, the comparison in
Section~\ref{sec:comparison} gives the lower bound
$\exp\bigl(c(\log N)^2/(\log\log N)^2\bigr)$ with an absolute $c>0$.
In each characteristic admitting ports, Theorem~\ref{thm:stretch}
gives $2^{N^\varepsilon}$ for some $\varepsilon>0$ depending on that
characteristic. This exceeds the first bound with the same absolute
$c$ once $N$ is sufficiently large in the given characteristic.
The two cases exhaust all characteristics, including zero. They give
the stated bound for $\HC_N$ over every field and for $\Per_N$ over
every field of characteristic different from two.
\end{proof}
